\begin{problem}[第34届物理决赛][电场磁场束缚粒子]{84}
    如图，某圆形薄片材料置于 $xOy$ 水平面上，圆心位于坐标原点 $O$；$xOy$ 平面上方存在大小为 $E$、沿 $z$ 轴负向的匀强电场，以该圆形材料为底的圆柱体区域内存在大小为 $B$、沿 $z$ 轴正向的匀强磁场，圆柱体区域外无磁场。从原点 $O$ 向 $xOy$ 平面上方的各方向均匀发射电荷量为 $q$、质量为 $m$、速度大小为 $v$ 的带正电荷的粒子。粒子所受重力的影响可忽略，不考虑粒子间的相互作用。
    \begin{subproblem}
        \begin{subsubproblem}
            若粒子每次与材料表面的碰撞为弹性碰撞，且被该电场和磁场束缚在上述圆柱体内的粒子占发射粒子总数的百分比为 $\eta = 50\%$，求该薄片材料的圆半径 $R$。
        \end{subsubproblem}
        \begin{subsubproblem}
            若在粒子每次与材料表面碰撞后的瞬间，速度竖直分量反向，水平分量方向不变，竖直方向的速度大小和水平方向的速度大小均按同比例减小，以至于动能减小 $10\%$。
            \begin{subsubsubproblem}
                \begin{subsubsubsubproblem}
                    求在粒子射出直至它第一次与材料表面发生碰撞的过程中，粒子在 $xOy$ 平面上的投影点走过路程的最大值；
                \end{subsubsubsubproblem}
                \begin{subsubsubsubproblem}
                    对（i）问中投影点走过路程最大的粒子，求该粒子从发射直至最终动能耗尽而沉积于材料表面的过程中走过的路程。
                \end{subsubsubsubproblem}
            \end{subsubsubproblem}
        \end{subsubproblem}
    \end{subproblem}
\end{problem}